% Part: normal-modal-logic
% Chapter: syntax-and-semantics
% Section: entailment

\documentclass[../../../include/open-logic-section]{subfiles}

\begin{document}

\olfileid{nml}{syn}{ent}

\olsection{অনুসিদ্ধান্ত}

\begin{explain}
  জগতে সত্যতার সংজ্ঞা থেকে !!{formula}s-এর মধ্যে অনুসিদ্ধান্ত-সম্পর্ক
  সংজ্ঞায়িত করা যায়। !!^a{formula}~$!B$ থেকে~$!A$ অনুসিদ্ধান্ত হয়
  তখনই এবং কেবল তখনই, যখন $!B$ সত্য হলেই $!A$-ও সত্য। এখানে
  ``হলেই'' কথাটি যে মডেলই বিবেচনা করি এবং সেই মডেলের যে জগৎই
  বিবেচনা করি—উভয় ক্ষেত্রেই প্রযোজ্য।
\end{explain}

\begin{defn}
  $\Gamma$ যদি !!{formula}s-এর একটি সেট এবং $!A$ !!a{formula} হয়,
  তবে $\Gamma$ থেকে~$!A$ \emph{অনুসিদ্ধান্ত হয়}, সংকেতে
  $\Gamma \Entails !A$, তখনই এবং কেবল তখনই, যখন প্রত্যেক মডেল
  $\mModel{M} = \tuple{W, R, V}$ এবং $w \in W$-এর জন্য,
  প্রত্যেক $!B \in \Gamma$-তে $\mSat{M}{!B}[w]$ হলে
  $\mSat{M}{!A}[w]$। $\Gamma$-তে একটিমাত্র !!{formula}~$!B$
  থাকলে লিখি $!B \Entails !A$।
\end{defn}

\begin{ex}
  !!a{formula} থেকে আরেকটি অনুসিদ্ধান্ত হয় দেখাতে,
  $\mSat{M}{}[w]$-এর সংজ্ঞা ব্যবহার করে সব মডেল সম্পর্কে যুক্তি দিতে
  হয়। যেমন, $p \lif \Diamond p \Entails \Box\lnot p \lif \lnot p$
  দেখাতে এভাবে যুক্তি দেওয়া যায়। যেকোনো মডেল
  $\mModel{M} = \tuple{W, R, V}$ এবং $w \in W$ নিয়ে ধরা যাক,
  $\mSat{M}{p \lif \Diamond p}[w]$। দেখাতে হবে,
  $\mSat{M}{\Box\lnot p \lif \lnot p}[w]$। উল্টোটি ধরে নিই।
  তাহলে $\mSat{M}{\Box\lnot p}[w]$ এবং
  $\mSat/{M}{\lnot p}[w]$। $\mSat/{M}{\lnot p}[w]$ থেকে
  $\mSat{M}{p}[w]$। অনুমান অনুসারে
  $\mSat{M}{p \lif \Diamond p}[w]$, তাই
  $\mSat{M}{\Diamond p}[w]$। $\mSat{M}{\Diamond p}[w]$-এর
  সংজ্ঞা থেকে এমন কোনো $w'$ আছে যাতে $Rww'$ এবং
  $\mSat{M}{p}[w']$। আবার $\mSat{M}{\Box \lnot p}[w]$
  হওয়ায় $\mSat{M}{\lnot p}[w']$—এটি স্ববিরোধ।

  !!a{formula}~$!B$ থেকে আরেকটি~$!A$ অনুসিদ্ধান্ত হয় না দেখাতে
  একটি বিপরীত উদাহরণ দিতে হয়। অর্থাৎ, এমন মডেল
  $\mModel{M} = \tuple{W, R, V}$ দেখাতে হয় যার কোনো $w \in W$-তে
  $\mSat{M}{!B}[w]$, অথচ $\mSat/{M}{!A}[w]$। দেখাই যে
  $p \lif \Diamond p \Entails/ \Box p \lif p$।
  \olref{fig:counterex}-এর মডেলটি বিবেচনা করো। সেখানে
  $\mSat{M}{\Diamond p}[w_1]$, ফলে
  $\mSat{M}{p \lif \Diamond p}[w_1]$। কিন্তু
  $\mSat{M}{\Box p}[w_1]$ হলেও $\mSat/{M}{p}[w_1]$;
  তাই $\mSat/{M}{\Box p \lif p}[w_1]$।
  \begin{figure}
  \begin{center}
    \begin{tikzpicture}[modal]
      \node[world] (w1) [label=right:\mFalse{p}]{$w_1$};
      \node[world] (w2) [label=right:\mTrue{p}, above left=of w1]{$w_2$};
      \node[world] (w3) [label=right:\mTrue{p}, above right=of w1] {$w_3$};
      \draw[->] (w1) to (w2);
      \draw[->] (w1) to (w3);
    \end{tikzpicture}
  \end{center}
  \caption{$p \lif \Diamond p
  \Entails/ \Box p \lif p$-এর বিপরীত উদাহরণ।} % BN-SRC-341
  \ollabel{fig:counterex}
  \end{figure}

  অনেক সময় আরও সরল বিপরীত উদাহরণই যথেষ্ট।
  $\mModel{M'} = \tuple{W', R', V'}$ মডেলে % BN-SRC-340
  $W' = \{w\}$, $R' = \emptyset$ এবং $V'(p) = \emptyset$ ধরলেও
  বিপরীত উদাহরণ পাওয়া যায়। $\mSat/{M'}{p}[w]$ বলে
  $\mSat{M'}{p \lif \Diamond p}[w]$। $w$ থেকে কোনো জগৎ
  অভিগম্য নয় বলে $\mSat{M'}{\Box p}[w]$; তাই
  $\mSat/{M'}{\Box p \lif p}[w]$।
\end{ex}

\begin{prob}
  দেখাও, $\Box (!A \land !B) \Entails \Box !A$।
\end{prob}

\begin{prob}
  দেখাও, $\Box (p \lif q) \Entails/ p \lif \Box q$ এবং
  $p \lif \Box q \Entails/\Box (p \lif q)$।
\end{prob}

\end{document}
